\begin{claim}{E.11}{Proved}{The intermediate-subcategory assertion under replacement data}
Let $\mathcal M$ have pullbacks, let $W$ contain isomorphisms and satisfy two-out-of-three, and suppose that every arrow factors as a $W$-map followed by an intrinsic $W$-fibration. Suppose in addition that for every base $S$ there is a functor
\[
R_S:\mathcal M/S\longrightarrow\operatorname{Fib}_W S
\]
and a natural transformation $\eta_S:1_{\mathcal M/S}\to i_SR_S$ whose components are in $W_S$, where $i_S$ denotes the full inclusion. Then XII's axioms (a), (b), and (c), in their exact formulations A.10, hold. In particular, for every full subcategory
$\operatorname{Fib}_W S\subseteq\mathcal H\subseteq\mathcal M/S$, both inclusions induce equivalences at the restricted $W_S$.
\end{claim}
\begin{proof}
The factorization argument of E.7, applied to a $W$-map between two intrinsic fibrations over $S$, proves axiom (b) as follows. Factor an arbitrary $T\to S$ as $T\xrightarrow{i}T_0\xrightarrow{p}S$ with $i\in W$ and $p$ intrinsic. The pullback of the map along $p$ is weak. Both comparison maps from pullback over $T$ to pullback over $T_0$ are weak, because the two structural fibrations preserve $i$. Two-out-of-three gives weakness after pullback along $T\to S$.

For $u\in W_S$, naturality gives $R_S(u)\eta_X=\eta_Yu$, so $R_S(u)\in W_S$. Hence $R_S$ descends to localizations. Its restriction to $\mathcal H$ has values in $\operatorname{Fib}_W S\subseteq\mathcal H$. Since $\mathcal H$ is full, every comparison $\eta_X:X\to R_SX$ for $X\in\mathcal H$ is a morphism inside $\mathcal H$. After localization these comparisons give natural isomorphisms between the identity and inclusion followed by $R_S$, both on $\mathcal H$ and on the fibrant subcategory. They exhibit the first inclusion as an equivalence. The same argument with $\mathcal M/S$ exhibits the composite inclusion as an equivalence, hence also the second. This proves (c). Finally (b) identifies the two restricted weak-equivalence classes on fibrant objects, so (c) gives (a). No compatibility of the chosen $R_S$ for distinct bases is required for this assertion.
\end{proof}

\begin{claim}{E.12}{Refuted}{A six-object counterexample to XIIc B062}
The implication $\mathrm{(a)}+\mathrm{(b)}\Rightarrow\mathrm{(c)}$ of A.10 is false, even for a finite lattice and a strongly saturated class $W$. Both premises can hold at every base while the conclusion fails at one base. The same example refutes the assertion with infinity-categorical localizations in place of ordinary localizations.
\end{claim}
\begin{proof}
Let $\mathcal M$ be the poset category on $\{0,1,2,3,4,5\}$ with Hasse covers
\[
 0<1,\quad 0<2,\quad 1<3,\quad 2<3,\quad
 1<4,\quad 3<5,\quad 4<5.
\]
There is exactly one morphism $a\to b$ when $a\leq b$, and no morphism otherwise; composition is the unique morphism $a\to c$. Take
\[
 W=\{1_x\mid 0\leq x\leq5\}\ \cup\
 \{1\to3,\ 1\to4,\ 2\to3\}.
\]
This specifies all objects, arrows, compositions and weak equivalences. The object $0$ is initial and $5$ terminal. Pullbacks are meets. For comparable elements the meet is the smaller element; the remaining unordered pairs are $\{1,2\}$, $\{2,4\}$ and $\{3,4\}$, whose meets are respectively $0,0,1$. Thus all pullbacks exist; the finite meet semilattice with maximum $5$ is a lattice.

Write $\mathcal F_s$ for the full subcategory of intrinsic fibrations over $s$. The following table gives every slice and every $\mathcal F_s$:
\[
\begin{array}{c|l|l}
s&\operatorname{Ob}(\mathcal M/s)&\operatorname{Ob}(\mathcal F_s)\\ \hline
0&0&0\\
1&0,1&0,1\\
2&0,2&0,2\\
3&0,1,2,3&0,3\\
4&0,1,4&0,1,4\\
5&0,1,2,3,4,5&0,3,5
\end{array}
\]
Indeed $x\leq s$ is intrinsic precisely when
\[
 a\wedge x\longrightarrow b\wedge x\ \in W
 \qquad\text{for every }(a\to b)\in W\text{ with }b\leq s.
\]
Over $3$, the candidate $1$ fails on $2\to3$, giving $0\to1$, and the candidate $2$ fails on $1\to3$, giving $0\to2$. Over $5$ these failures persist, and the candidate $4$ fails on $2\to3$, giving $0\to1$. For $x=3$ over $5$, the three nonidentity weak arrows pull back respectively to $1\to3$, $1_1$ and $2\to3$, all weak. Pullback along $0$ or along the identity of a base also preserves $W$. Over $4$ the only nonidentity weak arrow is $1\to4$; its pullbacks along $0,1,4$ are respectively $1_0,1_1,1\to4$. The slices over $0,1,2$ have only identity weak arrows. These observations prove every entry of the table.

We next check axiom (b) at every base. For $s\ne4$, the restriction of $W_s$ to $\mathcal F_s$ consists solely of identities, which are universal weak equivalences. For $s=4$, the additional arrow $1\to4$ is universally weak in that slice by the three pullbacks just computed. Therefore $W_s^u=W_s$ on $\mathcal F_s$ for all six bases.

For (a), the bases $0,1,2,4$ need no argument beyond (b), since $\mathcal F_s=\mathcal M/s$. At base $3$, define
\[
 R_3(0)=0,\qquad R_3(1)=R_3(2)=R_3(3)=3.
\]
This is a relative functor $\mathcal M/3\to\mathcal F_3$, is the identity on $\mathcal F_3$, and has a natural transformation $1_{\mathcal M/3}\to i_3R_3$ with weak components $1\to3$, $2\to3$, and identities. It therefore induces an inverse to the localized inclusion.

At base $5$, let $i:\mathcal F_5\hookrightarrow\mathcal M$ be the inclusion and define functors $R:\mathcal M\to\mathcal F_5$ and $T:\mathcal M\to\mathcal M$ by
\[
\begin{array}{c|rrrrrr}
x&0&1&2&3&4&5\\ \hline
R(x)&0&3&3&3&3&5\\
T(x)&0&1&2&3&1&5 .
\end{array}
\]
Monotonicity is checked on the seven displayed Hasse covers. Both functors preserve $W$: $R$ sends the three nonidentity weak arrows to $1_3$, whereas $T$ sends them to $1\to3$, $1_1$, and $2\to3$, respectively. Since $T(x)\leq x$ and $T(x)\leq R(x)$ for every $x$, there are natural transformations
\[
 1_{\mathcal M}\ \longleftarrow\ T\ \longrightarrow\ iR.
\]
All components are weak: the only nonidentity component of the left transformation is $1\to4$ at $x=4$; those of the right transformation are $1\to3$ at $x=1,4$ and $2\to3$ at $x=2$. In a poset every such naturality square commutes by uniqueness of its arrows. Also $Ri=1_{\mathcal F_5}$, and both transformations restrict to identities on $\mathcal F_5$. Localizing the two natural transformations gives a natural isomorphism $1\cong iR$. Together with $Ri=1$ this proves that the inclusion induces an equivalence. By (b), this is exactly the comparison in (a), including its prescribed class $W_5^u$. Thus (a) holds at every base.

For completeness, $W$ satisfies all the stated saturation requirements. Inspection of the arrows gives
\[
 (a\to b)\in W\quad\Longleftrightarrow\quad R(a)=R(b).
\]
The inverse image of identities under a functor to a poset satisfies two-out-of-three. Since $R$ inverts $W$, it factors through $\mathcal M[W^{-1}]$. Any original arrow becoming invertible there must map to an isomorphism in the chain $\mathcal F_5=\{0<3<5\}$, hence to an identity, and hence belongs to $W$. This proves strong saturation as well as the isomorphism requirements.

Finally take the full intermediate subcategory
\[
 \mathcal F_5=\{0,3,5\}\ \subseteq\
 \mathcal H=\{0,3,4,5\}\ \subseteq\ \mathcal M/5.
\]
It is the diamond with $0<3<5$ and $0<4<5$. The restriction $W_5\cap\operatorname{Mor}(\mathcal H)$ consists only of identities, as does its restriction to $\mathcal F_5$. Consequently the first localized inclusion is the ordinary inclusion of a three-object chain into this four-object diamond. It is not essentially surjective: $4$ is not isomorphic to any object of $\mathcal F_5$. The second localized inclusion is not full, since $\operatorname{Hom}_{\mathcal H}(3,4)$ is empty, whereas $3\leftarrow1\to4$ gives an isomorphism from $3$ to $4$ in $\mathcal M[W^{-1}]$. Both required equivalences in (c) therefore fail.

The enhanced assertion follows from the same certificate, without any assumption about discrete mapping spaces of a general localization. The relative functors and componentwise weak natural transformations above descend under infinity-categorical localization to functors and natural equivalences, exhibiting the localized inclusions in (a) as equivalences. The localizations of $\mathcal F_5$ and $\mathcal H$ at identity arrows are their nerves; the failure of essential surjectivity persists there. Thus the example also meets the stronger enhanced premises and violates the enhanced conclusion.
\end{proof}

The independent program \texttt{code/a\_b062\_verify.py} verifies every arrow, every composition, every meet, every intrinsic-fibration test, every universal-equivalence test, and the displayed functors and natural transformations. Its certificate does not use the search program or a bounded-path approximation to localization. The example does not have enough intrinsic fibrations: $4\to5$ cannot factor as a weak arrow followed by an intrinsic fibration, since the only weak arrow with source $4$ is its identity. Nor does it have a coaugmented replacement as in E.11. The weak natural zigzag at base $5$ passes through $T(4)=1$, an object excluded from $\mathcal H$; this is exactly why it proves (a) without proving (c).
